\documentclass[10pt,twocolumn]{article}
\usepackage[a4paper,margin=1.9cm,columnsep=0.7cm]{geometry}
\usepackage{mathptmx}
\usepackage[T1]{fontenc}
\usepackage{microtype}
\usepackage{graphicx}
\usepackage{booktabs}
\usepackage{caption}
\usepackage{authblk}
\usepackage[hidelinks]{hyperref}
\usepackage{xcolor}

\captionsetup{font=small,labelfont=bf,skip=5pt}
\raggedbottom
\setlength{\parindent}{1.1em}
\definecolor{mono}{HTML}{2F6F9F}
\definecolor{dia}{HTML}{C25A3C}

\title{\vspace{-1.2em}\bfseries Monologue or Dialogue?\\[2pt]
\large Measuring Conversational Structure in the Reasoning Traces of Large Language Models}
\author[ ]{Oskar Engl}
\affil[ ]{\small Johannes Kepler University Linz, Austria}
\affil[ ]{\small \textit{oskar.engl@ik.me}}
\date{\small Working version, July 2026}

\begin{document}
\twocolumn[
\begin{@twocolumnfalse}
\maketitle
\begin{abstract}
\noindent
A reasoning model such as DeepSeek-R1 produces a long stretch of thinking before it answers. For
text written by one model in one pass, it behaves oddly: questions receive answers, objections
receive concessions, and the argument turns against itself and starts again. Read as text, is a
reasoning trace a monologue or a dialogue? We put that question to three text-analysis methods, and
required each to prove itself first on four corpora whose voice structure is known: people thinking
aloud over a puzzle, groups deliberating, one Qwen-2.5-7B instance solving the same puzzle alone, and
four role-prompted instances debating it. Only a method that separates monologue from dialogue in all
four goes on to measure the traces. Two of the three do: one reads how far writing style shifts
between neighbouring sentences, the other how often a conversational move receives the reply it calls
for. Both place DeepSeek-R1's traces past the four-instance debate, and the two part ways once people
enter the comparison. The traces vary more than any dialogue we measured, real group deliberation
included, and complete question-and-answer pairs at a third the rate of a person thinking aloud
alone. Whether that comes from several voices, or from one voice writing conversationally, is not
something the text can settle.
\medskip
\end{abstract}
\end{@twocolumnfalse}
]

\begin{figure*}[t]
\centering
\includegraphics[width=0.96\textwidth]{fin_fig_grid_v2.pdf}
\caption{\textbf{Study design.} Four corpora of known voice structure (left) supply the validity
check: a method is carried forward only if it separates monologue from dialogue on all four known
cells. Two of the three methods do so; the document-level author count does not, and is set aside.
The carried-forward methods then measure the DeepSeek-R1 traces (\texttt{ai-reason}). The two AI
cells anchor the comparison, because they share both their producer and their task with the traces;
the two human cells first test whether a method can detect voice structure at all, and then serve as
reference points in the results.}
\label{fig:grid}
\end{figure*}

\section{Introduction}
Reasoning-oriented training has become standard in frontier models. Given a problem, a model such as
DeepSeek-R1 \cite{guo2025} first produces a long stretch of intermediate tokens
\cite{kambhampati2025}, commonly called a reasoning trace, and only then gives its final answer.
These traces have a notable feature: each is written by a single model in one pass, yet a trace often
contains questions and their answers, objections and concessions, and sudden reversals of direction.
In writing by people, features like these can mark a real exchange between participants, and they can
equally mark a single inner voice. This raises the question: is a reasoning trace's textual shape
closer to that of a monologue or of a dialogue?

Kim et al.~\cite{kim2026} argue that reasoning-trained models improve in part because they simulate
a ``society of thought,'' a set of internal perspectives that propose, criticise, and reconcile.
There are older parallels in cognitive science, which has understood reasoning as a social and
argumentative activity \cite{mercier2017} and the mind as a set of interacting agents
\cite{minsky1986}.

Our study rests on one assumption: if a trace is produced by several voices taking turns, the
structure of the writing should show it. Authorship analysis rests on the observation that writers
differ in small, largely involuntary habits, most robustly in how they use function words, and that
these habits are stable within an author and vary between authors \cite{stamatatos2009}. We
therefore treat the question in operational terms: a monologue should read as the work of one author
and a dialogue as the work of several, and we ask which of the two the traces are closer to.

\section{Design}
Reasoning-trained models have been widely available only since 2025 \cite{guo2025}, so there is
little prior work on the voice structure of their traces. Our approach is to compare against a grid
of corpora whose structure \emph{is} known (Figure~\ref{fig:grid}).

\textbf{Four known corpora, and the target.} We vary two properties across the grid: who produced
the text (a human or an AI) and whether it is a monologue or a dialogue. The four cells are human
monologue (think-aloud protocols, in which a person speaks their thoughts aloud while solving a task
\cite{wurgaft2025}; $N=600$), human dialogue (group deliberation from DeliData \cite{karadzhov2023};
$N=500$), AI monologue (a single Qwen-2.5-7B model solving a puzzle on its own; $N=600$), and AI
dialogue (four separate copies of Qwen-2.5-7B, each given a role, taking turns in a debate;
$N=600$). The fifth cell is DeepSeek-R1 itself: the model is given the same puzzle, and we take its
native thinking trace, the reasoning it produces before its final answer, as the document ($N=600$).
This is the cell whose structure we want to determine.

\textbf{The validity check.} Performance on the grid earns a method the right to judge the traces.
We define the check at the level of the cell and apply it to every method: does the method classify
each of the four known cells correctly as monologue or dialogue? A method is carried forward only if
it classifies all four correctly.

\textbf{Holding the task fixed.} The AI monologue, AI dialogue, and AI reasoning cells all solve the
same puzzle, Game of 24 \cite{yao2023}, a small arithmetic puzzle in which four given numbers must be
combined with the four basic operations to reach 24, and the human monologue cell uses the same
puzzle. This rules out a method separating the cells by their subject matter rather than by their
voice structure. The puzzle also sits near the limit of the Qwen model's ability, so the two AI
corpora contain genuine problem-solving rather than instant recall or uniform failure. The human
dialogue cell is the one exception, since it uses the Wason card task, another compact logic puzzle,
because no group-deliberation corpus exists for Game of 24. Producer and task are therefore constant
only across the three AI cells, so those two AI cells anchor the measurements and the human cells
stand beside them as reference points. We also removed the speaker labels from both dialogue corpora,
so that the corpora differ as little as possible apart from the number of speakers producing the
text.

\section{Methods}
The three methods read deliberately different signals, so that agreement among them carries more
weight than a single method alone. Each turns its signal into a monologue-or-dialogue call for a
cell.

\textbf{Method 1: counting authors across the document} \cite{nath2019}. This was the top-ranked
system in the PAN 2019 style change detection task, at $0.848$ accuracy on the binary single against
multi-author decision, where the next entry reached $0.604$ \cite{zangerle2019}. It divides a
document into paragraph-sized windows, describes each window by how often the 50 most common words
occur in it, and groups similar windows together. The number of groups is read as the number of
authors, which we reduce to a binary decision: one author (monologue) or more than one (dialogue). A
cell counts as correctly classified if most of its documents fall on the right side of that decision.

\textbf{Method 2: detecting style changes between sentences} \cite{schmidt2025}. This method
operates sentence by sentence, splitting documents with the Punkt tokenizer. For each pair of
neighbouring sentences it estimates the probability that the two were written by different authors,
using an encoder trained to capture writing style while ignoring topic \cite{patel2025}. Averaging
those probabilities over a document gives its raw score. That raw score drifts with document length
in our data, because longer documents run more words between one speaker change and the next, so a
fixed number of real switches is spread thinner. Each anchor cell therefore gets its own fitted curve
of raw score against length, and every document is measured against both curves at its own length.
The difference between those two distances is what we call the document's \emph{lean}: positive when
it sits nearer the dialogue curve, negative when it sits nearer the monologue curve. A cell is
classified by the sign of its mean lean.

\textbf{Method 3: measuring question-and-answer structure} \cite{schegloff1973}. This method reads
the text as a sequence of conversational moves. Tencent Hunyuan Hy3, a mixture-of-experts model with
295B parameters of which 21B are active per token, tags each utterance with its function using CAMS
\cite{duran2022}, a standard scheme whose 27 labels include tags such as
\texttt{propositionalQuestion}, \texttt{answer} and \texttt{confirm}. Three signals were specified in
advance from these tags, with one designated primary: how often an \emph{adjacency pair} is
completed. An adjacency pair is the basic sequential unit of conversation, a two-part exchange in
which the first part makes a specific second part due, such as a question followed by an answer, or a
proposal followed by its acceptance; it is the mechanism that binds one participant's turn to
another's rather than merely placing them side by side \cite{schegloff1973}. A single writer can
produce completed pairs as well, and think-aloud protocols do so often, so the rate carries
information only against a comparable baseline. The method is therefore validated by checking that
the dialogue documents complete pairs at a higher rate than their monologue counterpart within each
producer. The two remaining signals compare whole distributions of tags, and of transitions between
tags, across cells.

\section{Results}

\subsection{Which methods qualify}
Two of the three methods, Method~2 and Method~3, classify all four known cells correctly and are
carried forward to give their judgment on the traces. Method~1 misclassifies both dialogue cells and
is set aside.

\textbf{Method 1 fails the check.} On single-author text the method is nearly perfect, correctly
labelling $0.975$ of the human monologue documents and $0.998$ of the AI monologue documents. On
multi-author text it collapses: it recovers $0.423$ of the AI debates, so a majority of four-way
debates are read as the work of one author, and \emph{none} of the human debates, calling all $500$
of them single-author. It therefore places both known dialogue cells on the monologue side.

The human result has a likely explanation in the corpus. DeliData turns average $8.4$ words, with a
median of $6$, against $166$ words per turn in the AI debate. Method~1 describes paragraph-sized
windows by their function-word profile, and at six words per turn a single window absorbs several
speakers at once, leaving little for the clustering to separate. A method that reads a four-way
debate as one author, and that never flags the human debate as multi-voiced, has not shown that it
can detect voices, so it does not judge the traces. Its own reading of the traces is recorded in
Table~\ref{tab:verdict} for completeness.

\textbf{Deviation from the pre-registered gates.} This failure is also where our procedure departs
from what we fixed in advance. We originally pre-registered a separate statistical gate for each
method rather than the single cell-level check used here. Under those gates Method~1 cleared its
accuracy baseline and would have passed, while Method~3 was blocked by a label-reliability threshold
($\kappa \ge 0.70$, discussed below). The gates therefore admitted the method that reads a four-way
debate as the work of a single author, and excluded a method that separates all four known cells in
the right direction. They were an operationalisation of the requirement stated from the outset, that
a method separate monologue from dialogue on data of known structure before it is trusted on the
traces, and the cell-level check states that requirement directly, with the same test applied to all
three methods. We report against the cell-level check throughout, and say where a pre-registered gate
would have decided otherwise.

\textbf{Method 2 passes.} All four known cells fall on the correct side. The two monologue cells lean
negative, human monologue at $-0.0144$ and AI monologue at $-0.0066$; the two dialogue cells lean
positive, human debate at $+0.0036$ and AI debate at $+0.0200$. Pooled, the separation is strong,
with Cliff's $\delta=0.6211$, a rank-based measure of separation where $0$ is complete overlap and
$1$ no overlap. One qualification belongs with these numbers: the two AI cells define the curves the
lean is measured against, so their signs follow partly from construction, and the informative part of
the check is that the two human cells, which enter no fit, land on the correct sides.

\textbf{Method 3 passes, with a caveat about its labels.} Its completion rate separates dialogue from
monologue in the expected direction within each producer, $0.235$ against $0.095$ for the human cells
and $0.025$ against $0.003$ for the AI cells, with Cliff's $\delta=0.4717$.

Separating the cells and labelling every utterance correctly are two different achievements, and the
check establishes the first. We measured the tagger against a human annotator on $503$ utterances.
The two chose the same label in $0.742$ of cases, giving Cohen's $\kappa=0.54$, an agreement score
that corrects for chance matches, below the $0.70$ we had pre-registered. The disagreements
concentrate on one recurring decision rather than scattering across the scheme: $101$ of the $130$
involve \texttt{inform}, the general statement label, against a more specific label that also fits,
and they split in both directions, $56$ where the annotator chose the specific label and $45$ where
the tagger did. Unweighted $\kappa$ charges such a near miss as heavily as a disagreement between two
unrelated tags; the CAMS authors make this point themselves and recommend a measure that gives
partial credit \cite{duran2022}, under which agreement rises to a tree-weighted $0.62$. Reviewing the
disagreements also showed that part of them come from the single annotator applying the scheme less
strictly than its definitions require, which lowers $\kappa$ without implying that the tagger was
wrong. We carry Method~3 forward on the strength of its cell-level separation, and treat its labels
as the weakest link in its chain, with two further checks planned (Section~7).

\begin{figure*}[t]
\centering
\includegraphics[width=0.98\textwidth]{fin_fig_split_v3.pdf}
\caption{\textbf{All five corpora on each carried-forward method's own measure.} The shaded span is
the interval between the two AI anchors, the range this design calibrates directly. Both methods put
the traces past the AI debate. They part company on the comparison with people: in style switching
(A) the traces exceed every corpus measured, while in completed adjacency pairs (B) they sit close to
the AI floor and several times below human talk, including a person thinking aloud alone.}
\label{fig:split}
\end{figure*}

\begin{table}[t]
\centering
\small
\setlength{\tabcolsep}{4pt}
\caption{The three methods against the validity check, and what the two survivors report.
Method~1's own reading of the traces is bracketed, since it does not judge them.}
\label{tab:verdict}
\begin{tabular}{@{}lccl@{}}
\toprule
Method & Passes & Carried & Reading of the traces \\
\midrule
1: author count & no & no & [between the two poles] \\
2: style switch & yes & yes & past every corpus \\
3: adjacency pairs & yes & yes & past AI debate, far below human \\
\bottomrule
\end{tabular}
\end{table}

\subsection{What the two methods report}
The traces vary more than any dialogue we measured, and exchange less than people do even when
talking to themselves. Each method is read against its own measure in Figure~\ref{fig:split}, since
the two quantities do not share units.

On Method~2 the traces lean $+0.0433$, more than twice the AI debate's $+0.0200$, and further toward
dialogue than any other corpus in the study. The human debate is the surprise here, at $+0.0036$,
inside the interval between the two AI anchors. The reason is visible in the corpora: the AI debaters
were prompted into four distinct roles and write long turns, while DeliData participants write six
words at a time in a shared, flat register, so a sentence-level style detector finds far more
switching in the synthetic debate than in the real one.

On Method~3 the traces complete pairs at $0.035$, above the AI debate's $0.025$ and well above the AI
monologue's $0.003$. Set beside the human cells that lead shrinks: a person thinking aloud alone
reaches $0.095$ and human debate $0.235$, so on the measure that tracks conversational moves most
directly the traces sit at roughly a third of one person's self-talk and a seventh of real
conversation. The two remaining signals separate neither AI anchor from the other, which is itself
informative about where the dialogue signal lives: the AI monologue and the AI debate draw on much the
same mix of conversational moves, and what distinguishes them is how those moves are sequenced into
completed pairs.

\section{Discussion}
The traces vary more than any dialogue we measured, and exchange less than people do even when
talking to themselves. Everything behind that sentence is in the text: the methods read the writing a
model produced, rather than the computation that produced it.

\textbf{The scope of the claim.} A dialogic reading admits three explanations, and this design
separates only two of them. A trace may be self-talk, one voice posing questions and answering them.
It may be one producer composing an exchange deliberately. Or it may be several producers taking
turns. The grid anchors the first and the third: think-aloud is self-talk, and the human debate and
the four role-prompted instances are several producers. It holds no instance of the second, because
no corpus in the design was written by one producer composing an exchange in a single pass. The
measurements can therefore place the traces against self-talk and against multi-producer exchange,
and they cannot separate one producer composing an exchange from several producers holding one. A
novel is the everyday demonstration that those two can look alike on the page. Against the AI corpora
the result leans in the direction Kim et al.~\cite{kim2026} describe; against people it is plainly
distant from conversation on the measure that tracks conversational moves. We report the measurement
and leave it there.

\textbf{One voice can adopt a conversational register.} The competing account worth taking seriously
is that a single generator has learned to write in a self-questioning register. DeepSeek-R1 was
trained on outcome-based rewards, and its authors report that reflective re-evaluation emerged
without being targeted directly \cite{guo2025}. The two methods are exposed to that account
unequally. Method~3 is exposed squarely: a single writer poses questions and answers them, and
think-aloud does so at nearly four times the AI debate's rate, which is why its reading is a
comparison of rates, and why the traces' modest lead over the AI debate is weak evidence for
multiplicity. Method~2 is exposed more narrowly. It keys on features writers do not control
deliberately, so deliberate imitation is not an easy move, and it can establish that neighbouring
sentences read as the work of different authors. What it cannot establish is whether that comes from
several persistent voices or from one generator whose register moves as its approach moves. A change
of style accompanying a change of approach looks locally identical to a change of speaker, and
nothing in this design separates the two.

\textbf{What the AI debate anchor is worth.} Both methods place the traces past that anchor, so it
deserves scrutiny before the distance is taken at face value. Its four voices are four instances of
one model distinguished by role prompt, which limits how far apart their writing habits can drift,
and it also means a style-based method may separate that cell from the AI monologue partly on
role-driven register rather than on multiplicity itself. We take role differentiation to be part of
what a society-of-thought account posits, which makes that confound partly definitional. As a
dialogue exemplar the anchor is weak in absolute terms as well, with a completion rate of $0.025$
against the human debate's $0.235$. Passing it is a lower bar than approaching human conversation, so
the direction of the placement is the finding we defend rather than its size.

\textbf{How much weight each claim bears.} The two halves of the result rest on different footings.
That the traces sit on the dialogue side of the AI comparison rests on Method~2, which has the
cleanest validation of the three, and no method that passes the check places them on the monologue
side. That they fall short of human conversation rests on Method~3, whose labels are the weakest link
in the chain. The second half is the more interesting of the two and also the more fragile.

\section{Limitations}
The target is a single reasoning model, DeepSeek-R1, on a single family of puzzles, so the verdict
concerns this setting and does not yet generalise to reasoning traces at large.

The human monologue cell is transcribed speech while every other cell is written text, a difference
between cells that the design does not control. The two anchors also come from a different base model
than the target. The anchors share their model with each other, so the axis they define is not
distorted by that difference; what it can do is shift the traces along that axis, if DeepSeek-R1
simply writes differently from Qwen, and the size and direction of any such shift are unknown.

Document length affects Method~2. Most traces are longer than any document in the AI monologue cell,
so their measured distance from that anchor reaches past the lengths the anchor covers, and the size
of that gap is less certain than its direction. A separate $85$ documents across the calibration
cells were left out of Method~2 entirely: the Punkt tokenizer found fewer than two sentences in them,
mostly in the very short chat turns of the human debate, and a sentence-pair method needs at least
one pair to score.

\section{Future work}
The most direct next step is also the cheapest, and it addresses the ambiguity at the centre of the
study. The design lacks a cell written by one producer composing an exchange in a single pass, which
is what separates a deliberately dialogic monologue from a genuine dialogue (Section~5). One model
asked to write a four-way debate in a single output supplies that cell.

Two further extensions follow from the limitations above. Generating both anchors from the target
model itself, where the model permits a monologue and a role-prompted debate condition, would remove
the change of base model. And a recurrence test would sharpen what Method~2 can say: stable voices
taking turns should bring the same styles back across a document, where a single register drifting
with the approach need not. Method~1 attempted this with function-word clustering and lacked the
sensitivity; the same test built on sentence-level style embeddings is the natural successor.

On Method~3's labels, two checks are planned: a second independent coder with a step to resolve
disagreements, and re-tagging with a different model to measure how stable the tagger's labels are.

\begin{thebibliography}{99}
\small
\setlength{\itemsep}{0.6pt}\setlength{\parskip}{0pt}
\bibitem{guo2025} D.\ Guo et al. DeepSeek-R1 incentivizes reasoning in LLMs through reinforcement
learning. \textit{Nature} 645:633--638, 2025.
\bibitem{kambhampati2025} S.\ Kambhampati et al. Position: Stop anthropomorphizing intermediate
tokens as reasoning/thinking traces! \textit{arXiv:2504.09762}, 2025.
\bibitem{kim2026} J.\ Kim, S.\ Lai, N.\ Scherrer, B.\ Ag\"uera y Arcas, J.\ Evans. Reasoning models
generate societies of thought. \textit{arXiv:2601.10825}, 2026.
\bibitem{mercier2017} H.\ Mercier, D.\ Sperber. \textit{The Enigma of Reason.} Harvard Univ.\
Press, 2017.
\bibitem{minsky1986} M.\ Minsky. \textit{The Society of Mind.} Simon \& Schuster, 1986.
\bibitem{stamatatos2009} E.\ Stamatatos. A survey of modern authorship attribution methods.
\textit{JASIST} 60(3):538--556, 2009.
\bibitem{wurgaft2025} D.\ Wurgaft et al. Scaling up the think-aloud method. \textit{Proc.\ CogSci},
2025.
\bibitem{karadzhov2023} G.\ Karadzhov, T.\ Stafford, A.\ Vlachos. DeliData: a dataset for
deliberation in multi-party problem solving. \textit{Proc.\ ACM HCI} 7(CSCW2), 2023.
\bibitem{yao2023} S.\ Yao et al. Tree of thoughts: deliberate problem solving with LLMs.
\textit{NeurIPS} 36, 2023.
\bibitem{nath2019} S.\ Nath. Style change detection by threshold-based and window-merge clustering.
\textit{CLEF Working Notes}, 2019.
\bibitem{zangerle2019} E.\ Zangerle, M.\ Tschuggnall, G.\ Specht, B.\ Stein, M.\ Potthast.
Overview of the style change detection task at PAN 2019. \textit{CLEF Working Notes}, 2019.
\bibitem{schmidt2025} G.\ Schmidt et al. Style change detection using a sequential sentence-pair
classifier. \textit{CLEF Working Notes}, 2025. \textit{arXiv:2508.00675}.
\bibitem{patel2025} A.\ Patel et al. StyleDistance: content-independent style embeddings with
synthetic parallel examples. \textit{NAACL}, 2025.
\bibitem{duran2022} N.\ Duran, S.\ Battle, J.\ Smith. Inter-annotator agreement using the
Conversation Analysis Modelling Schema. \textit{Comm.\ Methods \& Measures} 16(3):182--214, 2022.
\bibitem{schegloff1973} E.\ Schegloff, H.\ Sacks. Opening up closings. \textit{Semiotica}
8(4):289--327, 1973.
\end{thebibliography}

\end{document}
